\section{Introduction}\label{sec:introduction}
Many Chebotarev counting problems are naturally observed through a loss of resolution: one may see only the image of a Frobenius class in a quotient of the Galois group, or only a restricted collection of quotient counts. The basic issue is then arithmetic and quantitative: which Artin spectral components remain visible after quotienting, and how much of the hidden exact-kernel spectrum can be reconstructed from the full quotient profile?

The main body answers this question for geometrically connected finite \etale{} abelian covers of smooth projective curves over finite fields, and then extends the formulas to smooth open curves by controlling the finitely many boundary points. Entropy and quotient language are used as quantitative observables for Frobenius distributions. The finite quotient model in the appendix is optional comparison material; it is logically independent of the Chebotarev theorem chain and is not used in any arithmetic proof.

\subsection{Arithmetic results}
Chebotarev density and its refinements relate prime-counting functions to the spectral data carried by Artin \(L\)-functions; see \cite{Tschebotareff1926,Serre1981,NeukirchSchmidtWingberg2008}. This paper uses the abstract spectral expansion only as a bookkeeping interface and then proves the concrete finite-field specialization needed for the main arithmetic theorem chain. The entropy and harmonic-analysis tools are classical \cite{CoverThomas2006,CsiszarKorner2011,DemboZeitouni1998,Diaconis1988,Terras1999}; the contribution is the exact-kernel recovery mechanism for quotient-observed Frobenius laws.

For a connected finite \etale{} abelian cover \(V\to U\) over \(\FF_q\), the degree-\(n\) primitive Frobenius classes define a law \(\mu_n\) on the Galois group \(G\).  For every quotient \(G/H\) we compare the pushforward of \(\mu_n\) with the uniform law and ask which character kernels are visible from the resulting entropy profile.  The normalized twisted amplitudes
\[
A_n(\chi)=nB_n(\chi)/q^{n/2}
\]
are expressed by the usual Artin reciprocal-zero formula, with the proper-divisor primitive contribution and the finitely many open-boundary Euler factors kept separate.  This separation is only a normalization of classical inputs: geometric Frobenius, the Artin Euler product, Weil bounds, Parseval, relative entropy, and subgroup-lattice M\"obius inversion are used in standard forms.

The main theorem package is Theorem \ref{thm:projective-finite-field-chebotarev-entropy-tomography} for projective curves and Theorem \ref{thm:finite-field-chebotarev-entropy-tomography} for open curves.  Uniformly in the quotient \(G/H\), the quotient \(\chi^2\)-defect is \(q^{-n}\) times the Artin energy visible on that quotient up to an error \(O(q^{-3n/2})\), and relative entropy is one half of the same expression.  M\"obius inversion on the subgroup lattice then recovers every exact-kernel energy
\[
E_n(K)=\sum_{\ker\chi=K}|A_n(\chi)|^2
\]
from the full quotient entropy profile, after the natural rescaling, with error \(O(q^{-n/2})\).

Theorem \ref{thm:finite-field-persistent-conductor-classification} gives a structural endpoint for the recovered energies.  If \(a(\chi)\) is the Artin conductor degree and \(d(\chi)=2g_C-2+a(\chi)\), purity and the Grothendieck--Ogg--Shafarevich formula show that \(d(\chi)\) is the total multiplicity of the unitary reciprocal-zero signal.  The quadratic imprimitive contribution has the same sign as the reciprocal-zero coefficients at the frequencies \(1\) and \(-1\), so cancellation is impossible.  A character contributes persistent second-order entropy exactly when \(d(\chi)>0\) or it is quadratic.  For a projective base this yields the sharp dichotomy that the coarsest persistent quotient is \(G\) in genus at least \(2\) and \(G/2G\) in genus \(1\).

The arithmetic contribution is therefore a classification of what quotient entropy
detects, not a new Chebotarev theorem or a new trace formula: the inputs are
standard purity and the Grothendieck--Ogg--Shafarevich formula, and what is proved
is which characters survive in the observable and which conductor data decides it.

Two further arithmetic features make the recovery usable.  First, the restricted-quotient problem has a sharp finite obstruction: Theorem \ref{thm:finite-field-exact-kernel-visibility-obstruction} identifies the subgroup-lattice rank condition for a chosen family of quotients to recover the same exact-kernel energies as the full profile.  Second, persistent secondary bias is treated without asserting pointwise phase recovery.  Corollary \ref{cor:cesaro-averaged-quotient-profile-persistent-support} and Theorem \ref{thm:kernel-chebotarev-persistent-energy-detection} recover the persistent exact-kernel support from Cesaro-averaged quotient profiles and from the unitary part of the finite-field Artin expansion.  The cyclic and biquadratic Kummer examples illustrate these arithmetic statements on explicit finite-field covers.

\begin{theorem}[Finite-field quotient entropy tomography]\label{thm:intro-main-arithmetic-tomography}
Let \(V\to U\) be a geometrically connected finite \etale{} abelian Galois cover of a smooth geometrically connected curve over \(\FF_q\), with group \(G\). For all sufficiently large \(n\), the family of quotient Frobenius laws \(\{\mu_n^H:H\le G\}\) determines, up to an error \(O(q^{-n/2})\) after the natural entropy rescaling, every exact-kernel Artin energy
\[
E_n(K)=\sum_{\substack{\chi\in\widehat G\setminus\{\ind\}\\ \ker\chi=K}}|A_n(\chi)|^2.
\]
More precisely, uniformly for all subgroups \(H,K\le G\),
\[
\chi^2\!\left(\mu_n^H\middle\|\mathrm{unif}_{G/H}\right)
=q^{-n}\sum_{\chi\in H^\perp\setminus\{\ind\}}|A_n(\chi)|^2+O(q^{-3n/2}),
\]
and
\[
E_n(K)=2q^n\sum_{J\supseteq K}\mu_{\mathrm{sub}}(K,J)
D\!\left(\mu_n^J\middle\|\mathrm{unif}_{G/J}\right)+O(q^{-n/2}).
\]
The fully quantified projective and open-curve versions, including the reciprocal-zero expression for \(A_n(\chi)\), are Theorem \ref{thm:projective-finite-field-chebotarev-entropy-tomography} and Theorem \ref{thm:finite-field-chebotarev-entropy-tomography}.
\end{theorem}

\begin{proof}
This is exactly the combination of the two cited arithmetic statements, with the displayed notation as defined in Definition \ref{def:finite-field-quotient-notation}. The first displayed formula is the uniform quotient \(\chi^2\)-asymptotic in Theorem \ref{thm:projective-finite-field-chebotarev-entropy-tomography} for projective \(U\) and in Theorem \ref{thm:finite-field-chebotarev-entropy-tomography} for open \(U\). The second displayed formula is the subgroup-lattice M\"obius inversion of the uniform entropy expansion in the same theorems. Proposition \ref{prop:finite-field-normalized-reciprocal-zero-formula} supplies the advertised expression of the normalized amplitudes in terms of Artin reciprocal zeros, proper-divisor terms, and boundary factors.
\end{proof}

At the quadratic level, supporting quotient identities explain the exactness of this recovery. For a surjection \(\pi:G_2\twoheadrightarrow G_1\), Proposition \ref{thm:kernel-chebotarev-quotient-relative-entropy-chain} gives the exact chain rule for relative entropy with respect to the uniform laws on \(G_2\), \(G_1\), and the fibers of \(\pi\). Proposition \ref{thm:kernel-chebotarev-visible-hidden-character-energy} is the corresponding \(L^2\)-statement: the \(\chi^2\)-defect splits orthogonally into a quotient-visible part, carried by characters factoring through \(\pi\), and a complementary invisible part. Theorem \ref{thm:kernel-chebotarev-quotient-lattice-tomography} packages the exact subgroup-lattice inversion. The later abstract second-order entropy statement is included only as supporting bookkeeping; the finite-field theorems just described are the proof path for the arithmetic article.

The final arithmetic statements isolate persistent secondary bias. The finite-field detection statements are Theorem \ref{thm:finite-field-exact-kernel-visibility-obstruction}, which proves the exact subgroup-lattice obstruction and recovery mechanism, Corollary \ref{cor:cesaro-averaged-quotient-profile-persistent-support}, which recovers persistent support directly from Cesaro-averaged rescaled quotient profiles, and Theorem \ref{thm:kernel-chebotarev-persistent-energy-detection}, which uses the recovered exact-kernel energies and the Cesaro-negligible finite-field expansion of Proposition \ref{prop:finite-field-unitary-artin-expansion} to identify the coarsest persistent quotient. The isolated-dominant-channel subset term and the pole/statistical convergence statements are auxiliary consequences of an abstract spectral model; they provide comparison language for dominant channels but are not prerequisites for the finite-field persistent-energy theorem.

\subsection{Appendix-level finite quotient model}
Appendices \ref{sec:groupoid-rigidity}--\ref{app:window6} contain a finite-dimensional fold-groupoid model. These results are not used to prove the Chebotarev theorems. They are included as a finite quotient-observable comparison: central idempotents replace character kernels, and Watatani-index moments replace quotient-visible quadratic energy. The appendix proves its own finite-dimensional claims, including fingerprint rigidity, central idempotents, and the window-\(6\) reconstruction calculation, without asserting an arithmetic identification with Chebotarev covers.

\subsection{Structure of the paper}
Section \ref{sec:entropy-chain} proves the finite-field quotient entropy tomography theorem and develops the exact quotient identities used in the recovery. Section \ref{sec:artin-pole} proves the finite-field exact-kernel obstruction, persistent-energy criterion, and conductor classification, and then records abstract dominant-channel consequences as optional comparison material. Appendices \ref{sec:groupoid-rigidity} and \ref{sec:central-idempotents} record the finite fold-groupoid comparison model and the auxiliary window-\(6\) computations.

\endinput
